\scanpage{35}
\uline{Inconvénient sérieux} de la \unsure{solution|notation}
$\unsure{\operatorname{ad}^{n}|\operatorname{ad}^{\circ}}(S,\unsure{J|I};\mathcal A)$ lorsque $\mathcal A$ satisfait $(\mathrm{AB}\,5)$, $\unsure{\mathfrak B|\mathfrak I}$ $S$ noethérien : les conditions de finitude naturelles (pour $\mathcal A$ loc \unsure{noeth}) sur des systèmes inductifs adiques, ou co. adiques, ou \unsure{fst ann. adiques}, ne semblent pas stables par \uline{conoyaux} [\unsure{dualemt}, dans $\unsure{\mathrm{cc}|\operatorname{ad}}(S,\unsure{J|I};\mathcal A)$ quand $\mathcal A$ satisfait $(\mathrm{AB}\,5^{\unsure{\alpha|*}})$, les \unsure{condns} de finitude artiniennes ne semblent pas stables par noyaux] !!

\deleted{\uline{Exemple} $A$ anneau [noeth], $f,g$ deux éléments de $A$, $M$ un objet de \deleted{\unsure{\(\mathcal A\)}} $\mathcal C$, \unsure{tel que} $(f,g)$ soit \unsure{$M$-régulier}. Soit
\[P=M_f/M=\varinjlim_n M/f^nM\]
($M/f^mM\longrightarrow M/f^nM$ est donné par $f^{n-m}$ ; le système inductif est à morphismes de transition injectifs). Soit
\[Q=\]
}
\mspage{34}
